\newpage
\deloppgave{Deloppgave 5.2 -- Studier og muligheter}
\setcounter{section}{0}

\section*{Introduksjon}
\addcontentsline{toc}{section}{Introduksjon}
Maskiningeniør er et bredt studium, og drift og vedlikehold er en sentral del av det. I løpet av dette semesteret har vi lært mye om hva drift og vedlikehold innebærer og hvordan undervisningen er lagt opp. Vi har også fått innblikk i arbeidshverdagen til en drift og vedlikeholdsingeniør og mulighetene etter bachelorstudiet. I denne refleksjonen diskuterer vi hva som motiverer oss til å velge studieretning drift og vedlikehold, og hvor viktig kombinasjonen av teori og praksis er for valget.

\section*{Motivasjon}
\addcontentsline{toc}{section}{Motivasjon}

Når man skal velge studieretning, er motivasjon en viktig faktor. Vår motivasjon for å velge drift og vedlikehold bygger på flere ulike inntrykk og erfaringer. Gjennom bedriftsbesøk og innsikt fra tidligere studenter har vi fått et bedre bilde av hva vedlikehold innebærer og hvilke muligheter studiet gir.

Det faglige innholdet i studieretningen er en viktig motivasjonsfaktor. Drift og vedlikehold kombinerer teknisk forståelse av maskiner og systemer med analytiske metoder som pålitelighetsberegning og tilstandsovervåking. Det er et fag der man må forstå hvordan ting fungerer, hvorfor de svikter og hva man kan gjøre for å forhindre det. Denne koblingen mellom teori og praktisk problemløsning er noe vi i gruppen synes er interessant. Før emnet hadde vi et nokså vagt bilde av hva en vedlikeholdsingeniør gjør. Nå ser vi at faget i stor grad handler om å ta beslutninger basert på data og analyse, ikke bare å skifte deler.

Vedlikehold åpner for mange muligheter i arbeidslivet og er etterspurt i flere bransjer, blant annet industri, energi, offshore og produksjon. For vår del virker dette særlig motiverende fordi studieretningen ikke låser oss til én bransje. Den samme kompetansen kan brukes i vannkraft, offshore, produksjon og samferdsel, men med ulike tekniske utfordringer. Fleksibiliteten gir også en trygghet i form av jobbsikkerhet, ettersom vedlikehold er noe det alltid vil være behov for. Det finnes jobbmuligheter både i Norge, fra sør til nord, og internasjonalt.

En annen faktor som motiverer er muligheten for en variert arbeidshverdag. Arbeidshverdagen kan variere mellom planlagt analysearbeid, akutt feilsøking og samarbeid med driftspersonell, noe som gjør fagfeltet mindre ensformig. For oss gjør denne variasjonen fagfeltet mer attraktivt enn en rolle med mer ensidige arbeidsoppgaver.

Studieretningen er godt tilrettelagt for læring. Klassene er relativt små, noe som gir tettere kontakt med foreleserne. RAMS-labben gir tilgang til relevant utstyr for praktiske øvinger, og flere av fagene legger opp til prosjektarbeid i samarbeid med lokal industri. Vi har sett eksempler på at studenter arbeider med reelle problemstillinger for lokale bedrifter. Dette er en fordel når man skal ut i jobb eller finne bacheloroppgave, fordi man allerede har kontakt med bedrifter og kan vise til resultater.

Etter fullført bachelor finnes det gode muligheter for videre studier på masternivå, både i Norge og i utlandet. Ved NTNU på Gløshaugen er det særlig masterprogrammet Maskinteknikk og produksjon (2-årig) som er relevant \parencite{ntnumiprod}. Innen dette programmet kan man velge studieretningen Industriell ledelse og systemfag, med fordypning innen sikkerhet, pålitelighet og vedlikehold. Denne masteren bygger videre på bachelorgraden og gir dypere innsikt i blant annet driftssikkerhet, vedlikeholdsstrategier og analyse av tekniske systemer.

\section*{Kombinasjon av teori og praksis}
\addcontentsline{toc}{section}{Kombinasjon av teori og praksis}
Studieretning drift og vedlikehold er fordelt mellom praktisk arbeid og teori. Labboppgaver gjennomføres i RAMS-labben, og prosjektoppgaver kan ved anledning gjøres i samarbeid med lokal industri.

Hvordan undervisningen er lagt opp har stor påvirkning på hvordan man lærer. Drift og vedlikehold har så langt hatt mye praktisk arbeid. Dette synes vi engasjerer mer enn ren teoretisk læring, fordi man kommer nærmere oppgavene man faktisk vil utføre i arbeidslivet.

Fagene videre i retningen bygger godt på hverandre \parencite{mast2003, mast2004, mast2006, mast2012}. MAST2003 gir grunnlag i kraftoverførende maskindeler og deres tendenser for svikt, mens MAST2004 bruker dette videre i en prosjektoppgave om en reell problemstilling innen drift og vedlikehold. At prosjektoppgaven ligner det man møter i næringslivet gjør den nyttig, ikke bare faglig, men fordi den gir oss en forsmak på hvordan vedlikeholdsproblemer faktisk løses i praksis. MAST2006 dekker analyse av levetid og driftssikkerhet, vedlikeholdsledelse og vedlikeholdsstyringssystemer, mens MAST2012 handler om prediktivt vedlikehold og bruk av tilstandsdata i modeller for prognose og diagnose. Det vi setter mest pris på med disse fagene er at læringsformen -- labøvinger, prosjektarbeid og interaktive forelesninger -- gjør at vi kan prøve ut metodene, ikke bare lese om dem.

\section*{Bacheloroppgaven}
\addcontentsline{toc}{section}{Bacheloroppgaven}
Bacheloroppgaven i sjette semester er en viktig del av studieretningen. Her får man mulighet til å jobbe med en større, reell problemstilling over lengre tid, ofte i samarbeid med en bedrift. Dette skiller seg fra de mindre prosjektoppgavene ved at man må jobbe mer selvstendig og ta ansvar for hele prosessen fra problemdefinisjon til løsningsforslag. Bacheloroppgaven fremstår derfor som et viktig punkt i studieløpet, fordi den gir mulighet til å teste om interessen for praktisk problemløsning og vedlikeholdsanalyse faktisk passer med arbeidsoppgavene i industrien.

\section*{Konklusjon}
\addcontentsline{toc}{section}{Konklusjon}
Drift og vedlikehold er en studieretning som gir kompetanse spisset mot industrien vi vil møte etter endt utdanning. Fagene har god variasjon mellom teori og praksis, noe som gir et solid læringsutbytte. Kombinasjonen av labbarbeid, prosjektoppgaver og bacheloroppgave i samarbeid med bedrifter gjør overgangen til arbeidslivet kortere. Gjennom semesteret har vi gått fra et vagt inntrykk av hva vedlikehold er, til en bedre forståelse av at det handler om systematisk analyse, risikostyring og datadrevne beslutninger. Det har gitt oss et godt grunnlag til å velge veien videre, enten det blir en mastergrad eller rett ut i jobb.
